\documentclass[computers,article,submit,moreauthors]{Definitions/mdpi}

\firstpage{1}
\makeatletter
\setcounter{page}{\@firstpage}
\makeatother
\pubvolume{1}
\issuenum{1}
\articlenumber{0}
\pubyear{2026}
\copyrightyear{2026}
\datereceived{}
\daterevised{}
\dateaccepted{}
\datepublished{}

\usepackage{amsmath}
\usepackage{amssymb}
\usepackage{booktabs}
\usepackage{multirow}
\usepackage{subcaption}
\usepackage{array}
\graphicspath{{./}}
\setlength{\headheight}{20pt}

\Title{SGBANDTI: A Controlled Evaluation of Rooted-Subgraph Drug Tokens with Bilinear Attention for Drug--Target Interaction Prediction}

\Author{Wen Li $^{1}$, Quan Long $^{1}$, Yanxiu Fu $^{1}$, Meng Li $^{1}$, Ju Xiang $^{2}$, Yanyan Yu $^{3}$ and Xiaojie Zhang $^{1,}$*}
\AuthorNames{Wen Li, Quan Long, Yanxiu Fu, Meng Li, Ju Xiang, Yanyan Yu and Xiaojie Zhang}

\address{%
$^{1}$ \quad School of Electronic Information, Hunan First Normal University, Changsha 410002, China; \\
$^{2}$ \quad School of Computer Science and Technology, Changsha University of Science and Technology, Changsha 410114, China; \\
$^{3}$ \quad Hunan Institute of Tuberculosis Prevention and Treatment, Hunan Provincial Chest Hospital Affiliated to Hunan Normal University, Changsha 410205, China}

\corres{Correspondence: xzhang6@gradcenter.cuny.edu}

\abstract{For drug--target interaction (DTI) prediction from molecular graphs and protein sequences, it remains unclear whether explicit rooted-subgraph extraction improves ranking beyond a standard GCN that retains the same atom-level tokens until cross-modal fusion. We developed SGBANDTI, which encodes atom-centered rooted subgraphs, applies convolutional layers to protein sequences, and fuses the features through masked low-rank bilinear association maps. Across five seeds on fixed random-pair splits, SGBANDTI achieved AUROCs of 0.9062 $\pm$ 0.0019 on BioSNAP and 0.9620 $\pm$ 0.0010 on BindingDB, slightly below DrugBAN but above the other baselines. A token-matched standard-GCN control achieved similar mean performance (AUROC 0.9051 $\pm$ 0.0043), indicating that the improvement over pooled whole-molecule variants was mainly associated with retaining atom-level tokens until fusion rather than with explicit rooted-subgraph extraction. On a canonical-SMILES identity-disjoint split SGBANDTI achieved the highest mean AUROC (0.8794 $\pm$ 0.0019), whereas performance decreased to 0.6345 $\pm$ 0.0182 on the protein-identity-disjoint split. Overall, SGBANDTI was competitive on random-pair and canonical-SMILES-disjoint settings but transferred poorly to exact-protein-sequence-disjoint targets.}

\keyword{drug--target interaction prediction; molecular subgraph; bilinear attention network; protein sequence modeling; BioSNAP; computational drug discovery}

\begin{document}

\section{Introduction}
Drug--target interactions (DTIs) connect bioactive compounds to the proteins that mediate therapeutic and adverse effects. Reliable DTI identification is therefore important for target discovery, compound prioritization, drug repurposing, and safety assessment \citep{yamanishi2008,bagherian2021}. Because biochemical assays and large-scale screening cannot exhaustively cover the expanding chemical and proteomic spaces, computational models are increasingly used to rank candidate pairs for subsequent investigation. The present study addresses binary DTI prediction, which estimates whether an interaction is represented in a benchmark, rather than affinity regression, which estimates a continuous binding value. This distinction determines the labels, loss function, model-selection procedure, and interpretation of predictions.

Early computational approaches integrated chemical and genomic similarities through kernel learning, semi-supervised inference, and matrix factorization \citep{xia2010,gonen2012}. More recent deep-learning methods learn task-specific representations directly from molecular structures and protein sequences \citep{wen2020}. Graph neural networks encode drugs by propagating information between bonded atoms \citep{duvenaud2015,kipf2017}; GNN-CPI and GraphDTA combine graph-based compound encoders with convolutional protein encoders; GraphDTA is an affinity-regression framework and is cited here as a related molecular-graph/protein-sequence representation rather than a binary-classification method \citep{tsubaki2019,nguyen2021graphdta}; and hierarchical models preserve atom-, motif-, or neighborhood-level chemical information before prediction \citep{liu2024higraphdti,zhang2021ngnn}. These developments improve molecular representation, but the stage at which local drug context is compressed can still determine which chemical environments remain available during drug--protein fusion.

Interaction modeling provides a complementary line of development. Concatenating modality-level vectors is efficient but cannot explicitly associate individual drug regions with protein-sequence regions. MolTrans constructs interactions between learned drug and protein substructures \citep{huang2021moltrans}, BACPI uses bidirectional attention \citep{li2022bacpi}, and DrugBAN and GraphBAN employ bilinear attention to model cross-modal feature relationships, with DrugBAN adding a domain-adaptation design and GraphBAN emphasizing inductive prediction \citep{bai2023drugban,hadipour2025graphban}. iNGNN-DTI further combines nested drug and structure-derived target graphs with cross-modal attention \citep{sun2024ingnndti}. These studies show that cross-modal interaction modeling is important, but they leave two questions unresolved for molecular-graph/protein-sequence DTI prediction. Existing comparisons do not clearly separate explicit rooted-subgraph extraction from the retention of atom-level drug tokens until cross-modal fusion. First (RQ1), does explicit rooted-subgraph extraction improve DTI ranking beyond a standard GCN that retains the same atom-level tokens until identical bilinear fusion? Second (RQ2), how does the full SGBANDTI configuration compare with representative baselines under random-pair, canonical-SMILES drug-disjoint, exact-protein-sequence target-disjoint, and BindingDB random-pair splits, and how well does held-out-pair performance reflect entity-disjoint generalization?

We therefore investigated whether retaining an atom-centered rooted subgraph for every valid drug position until low-rank bilinear fusion affects DTI ranking. Rooted neighborhoods preserve the local graph topology and chemical attributes surrounding each atom, which can distinguish atoms occupying different molecular environments before molecule-level compression. This hypothesis is deliberately representational rather than mechanistic: the model receives no protein structure, binding-site annotation, or atom--residue contact supervision, and its bilinear scores cannot be interpreted as physical contacts.

Neither the rooted-subgraph operator nor the bilinear module is new to this work; SGBANDTI's contribution is a task-specific combination of these components for molecular-graph/protein-sequence DTI ranking. To test this question, we developed SGBANDTI, which combines rooted $K$-hop molecular-subgraph encoding, sequential protein-sequence convolution, masked low-rank bilinear attention, and a multilayer classifier. The main contributions of this study are as follows:
\begin{itemize}
\item We design a token-matched standard-GCN control that separates explicit rooted-subgraph extraction from the retention of atom-level drug tokens until cross-modal fusion, and we find that the improvement over pooled whole-molecule variants was mainly associated with the latter rather than with explicit rooted-subgraph extraction.
\item We report a controlled negative result: explicit rooted-subgraph extraction did not yield a stable ranking gain over a standard GCN that retains the same atom-level tokens until identical bilinear fusion on the BioSNAP random split.
\item We reveal a clear asymmetry between drug-disjoint and target-disjoint generalization, identifying the latter as the main limitation of the architecture.
\item As a technical implementation, SGBANDTI couples atom-centered rooted-subgraph drug tokens with protein-sequence features through masked low-rank bilinear interaction, retaining per-atom neighborhood tokens until cross-modal aggregation.
\end{itemize}

\section{Materials and Methods}
\subsection{Overview of the SGBANDTI Framework}
Figure~\ref{fig:model_structure} presents the overall architecture of SGBANDTI. The model contains an input module, modality-specific encoders, a bilinear interaction block, and a multilayer perceptron (MLP) decoder. Drug SMILES strings are converted into molecular graphs, and target proteins are represented as embedded sequences. The drug encoder applies a nested subgraph strategy adapted from NGNN \citep{zhang2021ngnn}, whereas the target encoder applies a one-dimensional convolutional neural network (1D CNN). The resulting drug and target feature sets are then fused by a low-rank bilinear attention network, and the fused representation is mapped to a prediction score by the decoder. The protein cartoon in Figure~\ref{fig:model_structure} is illustrative only and is not used as an input feature.

\begin{figure}[H]
\centering
\includegraphics[width=0.95\textwidth]{model_structure.png}
\caption{Architecture of SGBANDTI. Drug SMILES strings are converted to molecular graphs and represented by rooted $K$-hop subgraphs. Target proteins are represented by integer-encoded amino-acid sequences and processed by three sequential one-dimensional convolutional layers. Valid drug and protein features are integrated through masked, unnormalized low-rank bilinear association maps, and an MLP produces the binary prediction score. Protein structures are illustrative only and are not supplied to the model. The implemented model does not use softmax normalization or max pooling in the bilinear module.}
\label{fig:model_structure}
\end{figure}

\subsection{Experimental and Technical Design}
The primary objective was to answer RQ1 and RQ2: whether explicit rooted-subgraph extraction improves ranking beyond a token-matched standard GCN, and how the resulting model compares with baselines under random-pair and entity-disjoint splits. BioSNAP was used for random interaction, unseen-drug, unseen-target, and architecture-replacement experiments, and BindingDB was used for cross-dataset replication. Random interaction, unseen-drug, and unseen-target scenarios were defined to distinguish held-out-pair performance from entity-disjoint generalization.

Five model-training runs used seeds 42, 52, 62, 72, and 82 on the same frozen split. The benchmark included conventional, graph-based, sequence-interaction, and nested-graph comparators, together with three replacement configurations and the full model, formed by replacing rooted-subgraph encoding, bilinear attention, or both, plus a token-matched standard-GCN control that uses the same atom featurizer, node count (290 including padding), drug mask, protein CNN, bilinear attention, decoder, and training protocol as the full model, but applies a standard multi-layer GCN over the whole molecular graph without explicit $K$-hop subgraph extraction; RQ1 is answered by comparing the full model with this control. AUROC and AUPRC were prespecified as primary threshold-independent metrics. For SGBANDTI, F1 score, sensitivity, specificity, and accuracy were secondary threshold-dependent metrics evaluated with a threshold selected exclusively from the validation set, and they were not used for the main between-model ranking. AUROC and average precision were computed with a shared evaluation script from archived per-sample probabilities for SGBANDTI and DrugBAN; per-sample probabilities are not archived for the remaining comparators. Comparative experiments used identical processed interaction files when the corresponding baseline implementation supported the shared protocol.

\subsection{Drug Feature Extraction with a \texorpdfstring{$K$}{K}-Hop Subgraph GNN}
Each drug molecule is represented as a graph $G = (V, E)$, where $V$ is the set of atoms and $E$ is the set of bonds. Only the bond-connectivity topology is consumed by the graph convolutions; explicit bond types and bond orders are not used as edge features. A stack of graph convolutional layers is applied inside the nested subgraph encoder to transform the original node attributes into hierarchical structural representations. At layer $l_0 + 1$, the representation of node $u$ is updated as
\begin{equation}
h_u^{(l_0+1)} = \sigma \left( \sum_{j \in \mathcal{N}(u)} \frac{1}{\sqrt{\hat{d}_u \hat{d}_j}} W^{(l_0)} h_j^{(l_0)} + b^{(l_0)} \right),
\end{equation}
where $\mathcal{N}(u)$ is the neighborhood of node $u$, $\hat{d}_u$ is the degree of node $u$ including a self-loop, $W^{(l_0)}$ and $b^{(l_0)}$ are the learnable weight matrix and bias, and $\sigma(\cdot)$ is a nonlinear activation function.

Instead of retaining only the final GCN layer, SGBANDTI concatenates the outputs from all GCN layers so that each node representation contains multi-scale neighborhood information:
\begin{equation}
\mathbf{z}_v = \mathbf{h}_v^{(1)} \mathbin{\|} \mathbf{h}_v^{(2)} \mathbin{\|} \cdots \mathbin{\|} \mathbf{h}_v^{(L)}.
\end{equation}
For each center node $u$, the $K$-hop subgraph $S_K(u)$ is summarized by average pooling over these concatenated node features:
\begin{equation}
\bar{\mathbf{z}}_u = \mathrm{AvgPool}\big(\{\mathbf{z}_v \mid v \in S_K(u)\}\big) = \frac{1}{|S_K(u)|}\sum_{v \in S_K(u)} \mathbf{z}_v,
\end{equation}
followed by a linear projection
\begin{equation}
\mathbf{g}_u = \mathbf{W}_p \bar{\mathbf{z}}_u + \mathbf{b}_p.
\end{equation}
This design retains one representation per centered subgraph and avoids whole-molecule pooling before interaction modeling.

In the implementation, SMILES strings were converted to bidirectional molecular graphs with DGL-LifeSci. Each real atom was encoded with the canonical 74-dimensional atom featurizer, and one padding-indicator bit was appended to obtain a 75-dimensional node feature. The raw node features were projected to a 128-dimensional embedding space, and the hidden representation of the drug encoder was also set to 128 dimensions. To support batched processing, molecular graphs were padded to a maximum of 290 nodes. The largest observed molecular graphs contained 288 atoms in BindingDB and 277 atoms in BioSNAP; molecules exceeding 290 atoms were outside the evaluated domain. Padding nodes carried zero-valued atom features plus an indicator bit set to 1, while real nodes used the same indicator bit with value 0. Padding nodes were disconnected from real atoms and had only self-loops. Real atoms also received self-loops once during graph construction. Therefore, padding nodes fixed tensor shape but did not serve as global information hubs. Formally, the feature of a real node is written as
\begin{equation}
\tilde{h}_i = [h_i; 0],
\end{equation}
whereas the feature of a virtual node is
\begin{equation}
\tilde{h}_v = [\mathbf{0}; 1],
\end{equation}
where $h_i$ denotes the original atom feature, $\mathbf{0}$ is a zero vector with the same feature dimension, and the last entry is the indicator bit specifying whether the node is virtual.

Rooted subgraphs were constructed for every graph position, including padding positions, but downstream masks ensured that only real atoms contributed to cross-modal aggregation. The main model used $K=2$. A shared three-layer GCN was applied independently to the rooted subgraphs; the outputs of the three layers were concatenated, mean pooled within each rooted subgraph, and linearly projected to 128 dimensions. This procedure yielded a fixed set of drug-position features and a drug mask derived from the padding-indicator bit.

\subsection{Target Protein Feature Extraction with a 1D CNN}
Target proteins were represented by embedded amino-acid sequences and processed by a sequential one-dimensional convolutional neural network. Protein sequences were converted to uppercase, and 25 standard or ambiguous amino-acid symbols were mapped to integer IDs 1--25. The value 0 was reserved for padding and unrecognized characters. Sequences longer than 1200 residues were truncated to the first 1200 N-terminal residues, discarding the remaining C-terminal portion, and shorter sequences were right padded to length 1200. This truncation affected 10.76\% of BindingDB records and 5.60\% of BioSNAP records; possible loss of C-terminal information is therefore a limitation of the fixed-length encoder.

Let $\mathbf{X} \in \mathbb{R}^{L \times d}$ denote the embedded target sequence, where $L$ is sequence length and $d$ is embedding dimension. At the $l_1 + 1$ convolutional layer, the feature map is computed as
\begin{equation}
\mathbf{h}^{(l_1+1)} = \sigma \left( \mathbf{h}^{(l_1)} * \mathbf{W}_c^{(l_1)} + \mathbf{b}^{(l_1)} \right),
\end{equation}
where $*$ denotes one-dimensional convolution, $\mathbf{W}_c^{(l_1)}$ is the convolution kernel, and $\mathbf{b}^{(l_1)}$ is the bias term. The embedding table contained 26 entries and mapped each token to a 128-dimensional vector, with index 0 used as the padding entry. The encoder contained three sequential convolutional layers with 128 output channels and kernel sizes of 3, 6, and 9, each followed by ReLU activation and batch normalization. Valid convolutions were used without zero padding; therefore, the final protein-feature length was
\begin{equation}
N_q = 1200 - (3-1) - (6-1) - (9-1) = 1185.
\end{equation}
Because valid convolutions were used without zero padding, the CNN reduced the protein length to 1185. The post-convolution protein mask was inferred from the nonzero integer residue length $\ell$, retaining only windows that are not fully padded:
\begin{equation}
m_j^{(q)} = \mathbb{I}\left(j < \min(\ell,1185)\right).
\end{equation}
Because unrecognized characters were encoded as 0, this mask is an approximation when such characters occur in the retained sequence.

\subsection{Low-Rank Bilinear Attention for Cross-Modal Fusion}
To preserve fine-grained interaction cues between drug substructures and local protein regions, SGBANDTI uses a bilinear attention network for pairwise cross-modal fusion. Let $\mathbb{V} \in \mathbb{R}^{N_v \times d_v}$ denote the drug features extracted by the subgraph GNN and $\mathbb{Q} \in \mathbb{R}^{N_q \times d_q}$ denote the target features extracted by the CNN. In the implemented model, $N_v=290$, $N_q=1185$, and $d_v=d_q=128$. The two modalities are first projected into a shared low-rank interaction space:
\begin{equation}
\mathbf{V}' = \mathcal{F}_v(\mathbf{V}), \quad \mathbf{Q}' = \mathcal{F}_q(\mathbf{Q}),
\end{equation}
where $\mathcal{F}_v$ and $\mathcal{F}_q$ are fully connected mappings in which dropout is applied before a weight-normalized linear projection followed by a ReLU activation. The bilinear hidden dimension was 256, the low-rank factor was $r=3$, and the expanded interaction dimension was therefore $256 \times r = 768$.

For the $s$-th attention head, the pairwise interaction score between the $i$-th drug feature and the $j$-th target feature is defined as
\begin{equation}
A_{s,i,j} = b_s + \mathbf{w}_s^\top \left(\mathbf{V}'_i \odot \mathbf{Q}'_j\right),
\end{equation}
where $\mathbf{w}_s$ and $b_s$ are learnable parameters. Drug and protein masks are applied before aggregation:
\begin{equation}
\widetilde{A}_{s,i,j} = A_{s,i,j} m_i^{(d)} m_j^{(q)}.
\end{equation}
Each head then aggregates pairwise responses through bilinear pooling:
\begin{equation}
\mathbf{z}_s = \sum_{i=1}^{N_v}\sum_{j=1}^{N_q} \widetilde{A}_{s,i,j}\left(\mathbf{V}'_i \odot \mathbf{Q}'_j\right),
\end{equation}
where $\odot$ denotes element-wise interaction in the shared feature space. To reduce redundancy introduced by low-rank expansion, grouped sum pooling is further applied:
\begin{equation}
\hat{\mathbf{z}}_s = \mathrm{AvgPool}_{r}(\mathbf{z}_s)\times r,
\end{equation}
which is implemented as average pooling over disjoint groups scaled by $r$ and therefore equals grouped sum pooling.
The final fused representation is
\begin{equation}
\mathbf{f} = \mathrm{BatchNorm}\left(\sum_{s=1}^{S}\hat{\mathbf{z}}_s\right).
\end{equation}
The implemented model used two attention heads, dropout 0.2 in the projection mappings, and no softmax normalization over the association maps. The masks $m^{(d)}$ and $m^{(q)}$ exclude padded drug positions and protein positions whose convolutional windows are entirely padding, so that only non-fully-padded protein windows contribute to the association maps. Thus, $A_{s,i,j}$ is an unnormalized latent feature-association score optimized for binary classification. It should not be interpreted as a probability, binding-site annotation, physical contact, or interaction energy.

\subsection{Decoder and Optimization}
The bilinear interaction feature $\mathbf{f}$ is processed by a four-layer MLP decoder with dimensional path $256 \rightarrow 512 \rightarrow 512 \rightarrow 128 \rightarrow 1$. Linear transformation, ReLU activation, and batch normalization were applied sequentially after each of the first three layers. The final output is passed through a sigmoid function to obtain the prediction score:
\begin{equation}
p = \mathrm{Sigmoid}(W\mathbf{f} + b).
\end{equation}

The model is trained by minimizing binary cross-entropy over the observed interaction pairs:
\begin{equation}
\mathcal{L}(\theta) = - \frac{1}{|\mathcal{B}|} \sum_{(d_i, t_j, y_{ij}) \in \mathcal{B}}
\Bigl[ y_{ij} \log \hat{y}_{ij} + (1 - y_{ij}) \log (1 - \hat{y}_{ij}) \Bigr].
\end{equation}
Because the decoder applied a sigmoid transformation in its final layer, binary cross-entropy was computed directly from the resulting probabilities using \texttt{BCELoss}; no additional sigmoid transformation was applied inside the loss function. In the reported experiments, the maximum number of epochs was 150, the batch size was 64, the learning rate was $5 \times 10^{-5}$, weight decay was 0, and the number of data-loading worker processes was 0. Training loaders were shuffled and dropped the final incomplete batch, whereas validation and test loaders preserved the input order and retained all samples.

\subsection{Datasets, Splits, and Evaluation Setting}
Let $\mathcal{D} = \{(d_i, t_j, y_{ij})\}$ denote the processed benchmark, where $d_i$ is a drug SMILES string, $t_j$ a protein amino-acid sequence, and $y_{ij} \in \{0,1\}$ the binary interaction label. The model learns a scoring function $f(d_i, t_j)$ estimating $P(y_{ij}=1)$, and AUROC and AUPRC are computed from the resulting scores on held-out pairs. This computational study evaluated binary DTI prediction on two public benchmark resources, BindingDB \citep{liu2007bindingdb} and BioSNAP \citep{zitnik2018}. The experiments used frozen processed benchmark files rather than reconstructing the datasets de novo from the complete source databases. Each processed record contained a drug represented by a SMILES string, a target represented by an amino-acid sequence, and a binary label.

The binary labels were interpreted as benchmark labels, not as complete biochemical truth. A label of 1 denotes a drug--target pair recorded as positive in the processed benchmark file. A label of 0 denotes a negative or unobserved pair supplied by the benchmark construction procedure. Such pairs should not be interpreted as experimentally confirmed non-binding pairs unless the source dataset explicitly documents experimental non-binding evidence. This distinction is especially important for DTI benchmarks, because many apparent negatives may represent untested or unreported interactions rather than measured absence of binding.

BindingDB curates experimentally reported drug--protein binding records, including continuous affinity measurements in the source database. The processed BindingDB files used here follow the binary benchmark setting adopted in prior DTI studies, including DrugBAN \citep{bai2023drugban}, and contain only binary labels. Provenance and traceability details are given in the Data Availability statement. Therefore, the BindingDB experiments should be interpreted as performance on the released binary benchmark representation, not as an independently reproducible reconstruction from raw BindingDB affinity records.

The BioSNAP benchmark contains reported positive associations and unlabeled pairs that are treated as negative examples in the distributed binary-classification files. The processed BioSNAP setting is consistent with recent binary DTI evaluation protocols, including iNGNN-DTI \citep{sun2024ingnndti}. Provenance and traceability details are given in the Data Availability statement. Therefore, BioSNAP 0 labels are treated operationally as sampled or unobserved negatives for benchmark discrimination, not as experimentally established non-interactions.

Protein sequences were converted to uppercase. Drug SMILES strings were parsed with RDKit and converted to canonical SMILES for pair-level duplicate detection. Records with unparsable SMILES strings were removed. Duplicate records were identified using the canonical drug representation and uppercase protein sequence; when duplicate pairs carried conflicting labels, the preprocessing script retained the positive label.

For random-split evaluation, the deduplicated pairs were shuffled with seed 42 and divided into training, validation, and test subsets at a ratio of 7:1:2 without class-stratified sampling. Identical processed drug--target pairs did not occur in more than one subset, although the same drug or target could occur in multiple subsets. This protocol therefore measured generalization to held-out pairs rather than to unseen entities. The split files were fixed across model runs; only parameter initialization, mini-batch order, and other stochastic training operations varied. Because the same drug or protein can recur across subsets, random-pair performance can overstate entity-level generalization relative to disjoint splits \citep{torrisi2022}, and five training seeds on one fixed partition quantify training stochasticity rather than uncertainty across data partitions.

The code also provides \texttt{unseen\_drug} and \texttt{unseen\_target} modes for entity-disjoint evaluation. A valid drug cold-start split requires every drug identity in the test set to be absent from training, whereas a target cold-start split requires every test protein to be absent from training. The unseen-drug split groups records by RDKit-canonical SMILES identity, preventing canonically equivalent drug strings from being assigned to both training and held-out subsets. The unseen-target split groups records by protein sequence identity. A canonical-SMILES-disjoint split is not scaffold-disjoint, and an exact-protein-sequence-disjoint split is not protein-family-disjoint: drugs with different canonical strings, or proteins with highly similar but non-identical sequences, can still span the training and test subsets.

\begin{table}[H]
\caption{Frozen split statistics used for the formal experiments. Sample counts and positives/negatives are reported per subset after preprocessing.}
\label{tab:dataset}
\scriptsize
\begin{tabular*}{\textwidth}{@{\extracolsep{\fill}}llllrrr@{}}
\toprule
\textbf{Dataset} & \textbf{Split} & \textbf{Set} & \textbf{Samples} & \textbf{Positive} & \textbf{Negative} \\
\midrule
\multirow{3}{*}{BioSNAP} & \multirow{3}{*}{Random} & Train & 19,220 & 9,725 & 9,495 \\
& & Validation & 2,746 & 1,370 & 1,376 \\
& & Test & 5,491 & 2,735 & 2,756 \\
\midrule
\multirow{3}{*}{BioSNAP} & \multirow{3}{*}{Unseen drug} & Train & 19,372 & 9,832 & 9,540 \\
& & Validation & 2,795 & 1,441 & 1,354 \\
& & Test & 5,290 & 2,557 & 2,733 \\
\midrule
\multirow{3}{*}{BioSNAP} & \multirow{3}{*}{Unseen target} & Train & 19,980 & 10,445 & 9,535 \\
& & Validation & 2,574 & 1,184 & 1,390 \\
& & Test & 4,903 & 2,201 & 2,702 \\
\midrule
\multirow{3}{*}{BindingDB} & \multirow{3}{*}{Random} & Train & 34,439 & 14,583 & 19,856 \\
& & Validation & 4,920 & 2,017 & 2,903 \\
& & Test & 9,840 & 4,074 & 5,766 \\
\bottomrule
\end{tabular*}
\noindent{\footnotesize{\textit{Note:} Random splits evaluate held-out drug--target pairs. Unseen-drug and unseen-target splits group held-out samples by canonical SMILES identity and protein sequence identity, respectively. BioSNAP contains 4,505 unique drugs and 2,181 unique proteins, and BindingDB contains 14,643 unique drugs and 2,623 unique proteins.}}
\end{table}

\subsection{Implementation Details, Model Selection, and Baselines}
Drug molecules were encoded with the canonical atom featurizer in DGL-LifeSci. Protein sequences were integer encoded, truncated to length 1200 when needed, and processed by the three-layer 1D CNN. Across five runs, predefined train/validation/test splits were fixed and only the random seed changed (42, 52, 62, 72, and 82). The seeds were applied to Python, NumPy, and PyTorch random-number generators. No early stopping criterion was used; all runs were trained for the maximum epoch budget, and the checkpoint used for testing was selected by validation AUROC:
\begin{equation}
e^* = \arg\max_e \mathrm{AUROC}_{\mathrm{val}}^{(e)}.
\end{equation}
Threshold-dependent metrics used a validation-selected threshold from the selected epoch:
\begin{equation}
\tau^* = \arg\max_\tau F_{1,\mathrm{val}}^{(e^*)}(\tau).
\end{equation}
The same $\tau^*$ was then applied once to test-set probabilities. Test labels were not used for checkpoint selection or threshold selection.

Table~\ref{tab:hyperparameter_config} lists the main hyperparameters. The SGBANDTI model contains approximately 1.07 million trainable parameters. On an NVIDIA GeForce RTX 4060 Laptop GPU with a batch size of 64 over the BioSNAP random training split, a single training epoch required approximately 55 s with 2.45 GB peak memory for SGBANDTI, 39 s with 2.05 GB for the pooled whole-molecule no-subgraph configuration, and 62 s with 2.46 GB for DrugBAN. Inference processed approximately 1,170 test samples per second for SGBANDTI (from cached subgraphs), 2,011 for the no-subgraph configuration, and 318 for DrugBAN, whose figure includes per-batch molecular-graph construction. Under a forward-only comparison with inputs pre-built (excluding graph construction), SGBANDTI processed approximately 1,092 samples/s (from cached subgraphs) and DrugBAN 1,301 samples/s; DrugBAN additionally required about 221 ms to construct each batch of molecular graphs. Retaining atom-centered subgraphs added about 16 s per epoch relative to whole-molecule pooling. We compared SGBANDTI with representative baselines spanning conventional machine learning, graph-based DTI prediction, sequence-interaction modeling, and nested graph representation: Random Forest (RF) \citep{breiman2001}, GNN-CPI \citep{tsubaki2019}, TransformerCPI \citep{chen2020}, MolTrans \citep{huang2021moltrans}, MGNDTI \citep{peng2024}, iNGNN-DTI \citep{sun2024ingnndti}, and DrugBAN \citep{bai2023drugban}. Baseline code versions, input modalities, seeds, and evaluation protocols are listed in the baseline protocol table in the Supplementary Materials. GraphBAN was not included in the main ranking table because the archived run used a transductive independent-graph setting that was not directly comparable with the shared interaction-split protocol.

\begin{table}[H]
\caption{Hyperparameter configuration used for SGBANDTI evaluation.}
\label{tab:hyperparameter_config}
\begin{tabular*}{\textwidth}{@{\extracolsep{\fill}}ll@{}}
\toprule
\textbf{Parameter} & \textbf{Value} \\
\midrule
Atom feature dimension & 75 \\
Drug node dimension & 128 \\
Drug graph size & 290 \\
Protein sequence length & 1200 \\
Protein dimension & 128 \\
1D CNN filters & [128, 128, 128] \\
1D CNN kernel sizes & [3, 6, 9] \\
Bilinear attention heads & 2 \\
Bilinear hidden dimension & 256 \\
Low-rank factor $r$ & 3 \\
MLP dimensions & 256, 512, 512, 128, 1 \\
Optimizer & Adam \\
Batch size & 64 \\
Learning rate & $5 \times 10^{-5}$ \\
Weight decay & 0 \\
Maximum epochs & 150 \\
Training seeds & 42, 52, 62, 72, 82 \\
Checkpoint selection & Highest validation AUROC \\
Threshold selection & Highest validation F1 at selected epoch \\
\bottomrule
\end{tabular*}
\end{table}

Hyperparameter selection. Architecture-scale settings, including the subgraph radius $K=2$, two bilinear attention heads, a hidden dimension of 128, a learning rate of $5 \times 10^{-5}$, a batch size of 64, and a maximum epoch budget of 150, follow the settings of the nested-graph and bilinear-attention baselines on which SGBANDTI builds. A maximum epoch budget of 150 was used for the main analysis, consistent with the frozen protocol for the reported results; the test set was not used for any hyperparameter choice.

\subsection{Statistical Analysis and Reproducibility}
For repeated-seed experiments, each metric is reported as the mean and sample standard deviation over five training seeds. Because all seeds used the same frozen data partition, these standard deviations quantify training stochasticity rather than uncertainty from split construction. Threshold-independent metrics, AUROC and AUPRC, are emphasized for between-model ranking because they are less sensitive to probability calibration and threshold shift. Paired bootstrap comparisons between SGBANDTI and DrugBAN were performed per training seed on shared test samples (1,000 resamples). Because test pairs can share drugs or proteins, these intervals do not strictly satisfy independent and identically distributed sampling, and the reported $\Delta$AUROC values are interpreted descriptively.

The software environment used Python 3.10.20, PyTorch 2.1.0, DGL 2.2.1 with CUDA 12.1 support, DGL-LifeSci 0.3.2, torch-scatter 2.1.2, NumPy 1.26.4, pandas 2.3.3, scikit-learn 1.6.1, and RDKit 2026.03.5, running on Windows 11. Training and evaluation used an NVIDIA GeForce RTX 4060 Laptop GPU (8 GB) with CUDA 12.1. Reproducibility from the processed benchmarks requires the frozen split files, per-seed prediction outputs, preprocessing records, figure-generation scripts, software versions, hardware notes, repository tag or commit identifier, and dataset access dates.

\subsection{Use of Generative AI in Writing}
During the preparation of this manuscript, the authors used OpenAI ChatGPT (web interface; accessed August--September 2026) and OpenAI Codex (web interface; accessed August--September 2026) to support language editing, structural revision, consistency checking, and LaTeX formatting. These tools were used for drafting and editing text only. The authors reviewed and edited all AI-assisted output and take full responsibility for the content. No generative AI tool was used to generate experimental data, design or run experiments, alter reported results, or fabricate references.

\section{Results}
\subsection{Random-Split Performance on BioSNAP and BindingDB}
The random-split experiments tested whether SGBANDTI remained competitive when drug--target pairs, but not necessarily drug or protein identities, were held out. All reported results were obtained from the frozen processed splits and archived per-seed outputs used for the final manuscript analyses. Frozen splits were available for BioSNAP random, BioSNAP unseen-drug, BioSNAP unseen-target, and BindingDB random experiments. All reported model summaries are means and sample standard deviations over seeds 42, 52, 62, 72, and 82 unless otherwise stated.

On the BioSNAP random split, SGBANDTI achieved an AUROC of $0.9062 \pm 0.0019$ and an AUPRC of $0.9132 \pm 0.0043$. DrugBAN achieved the highest mean AUROC and AUPRC on this split, exceeding SGBANDTI by 0.0038 AUROC and 0.0040 AUPRC. SGBANDTI nevertheless ranked above the remaining comparator methods, including MGNDTI, MolTrans, iNGNN-DTI, TransformerCPI, RF, and GNN-CPI.

The BindingDB random split showed the same pattern at a higher absolute performance level. SGBANDTI achieved an AUROC of $0.9620 \pm 0.0010$ and an AUPRC of $0.9482 \pm 0.0017$. DrugBAN again ranked first, with a small AUROC advantage of 0.0012. Thus, the random-split evidence supports SGBANDTI as an effective molecular-graph/protein-sequence DTI model, but it does not support a broad claim of dominance over DrugBAN.

This near-tie should be interpreted in light of the split design. In random pair-level partitions, the same drugs and proteins can appear in both training and test sets, so strong bilinear baselines can already learn shared entity-level and pair-neighborhood regularities. The small differences indicate similar mean ranking performance under the two fixed random-pair partitions, without establishing statistical equivalence between the methods. Paired bootstrap over the five seeds (1,000 resamples on shared test samples) confirmed that the SGBANDTI--DrugBAN differences were small: mean $\Delta$AUROC (SGBANDTI minus DrugBAN) was $-0.0039$ on BioSNAP random and $-0.0012$ on BindingDB, with 95\% confidence intervals crossing zero in most seeds. Because drugs and proteins recur across partitions, these experiments do not determine whether either architecture transfers to novel entities.

\begin{table}[H]
\caption{Random-split predictive performance on BioSNAP and BindingDB. Values are reported as mean $\pm$ sample standard deviation over five random seeds. The best value for each dataset and metric is shown in bold.}
\label{tab:model_performance}
\scriptsize
\setlength{\tabcolsep}{4pt}
\renewcommand{\arraystretch}{1.05}
\begin{tabular*}{\textwidth}{@{\extracolsep{\fill}}llcc@{}}
\toprule
\textbf{Dataset} & \textbf{Method} & \textbf{AUROC} & \textbf{AUPRC} \\
\midrule
\multirow{8}{*}{BioSNAP}
& DrugBAN        & \textbf{0.9100 $\pm$ 0.0023} & \textbf{0.9172 $\pm$ 0.0031} \\
& SGBANDTI       & 0.9062 $\pm$ 0.0019 & 0.9132 $\pm$ 0.0043 \\
& MGNDTI         & 0.8985 $\pm$ 0.0023 & 0.9046 $\pm$ 0.0024 \\
& MolTrans       & 0.8853 $\pm$ 0.0027 & 0.8918 $\pm$ 0.0034 \\
& iNGNN-DTI      & 0.8720 $\pm$ 0.0007 & 0.8771 $\pm$ 0.0013 \\
& RF             & 0.8402 $\pm$ 0.0008 & 0.8678 $\pm$ 0.0007 \\
& TransformerCPI & 0.8265 $\pm$ 0.0174 & 0.8456 $\pm$ 0.0108 \\
& GNN-CPI        & 0.7088 $\pm$ 0.0036 & 0.7229 $\pm$ 0.0010 \\
\midrule
\multirow{8}{*}{BindingDB}
& DrugBAN        & \textbf{0.9632 $\pm$ 0.0012} & \textbf{0.9516 $\pm$ 0.0019} \\
& SGBANDTI       & 0.9620 $\pm$ 0.0010 & 0.9482 $\pm$ 0.0017 \\
& MGNDTI         & 0.9500 $\pm$ 0.0014 & 0.9326 $\pm$ 0.0021 \\
& RF             & 0.9407 $\pm$ 0.0004 & 0.9209 $\pm$ 0.0005 \\
& MolTrans       & 0.9338 $\pm$ 0.0017 & 0.9086 $\pm$ 0.0054 \\
& iNGNN-DTI      & 0.9228 $\pm$ 0.0028 & 0.8951 $\pm$ 0.0024 \\
& TransformerCPI & 0.8921 $\pm$ 0.0012 & 0.8556 $\pm$ 0.0016 \\
& GNN-CPI        & 0.6548 $\pm$ 0.1712 & 0.5782 $\pm$ 0.2045 \\
\bottomrule
\end{tabular*}
\noindent{\footnotesize{\textit{Note:} Values are mean $\pm$ sample SD across five seeds. AUROC and AUPRC are threshold-independent metrics and were used as the primary between-method evidence. GraphBAN is excluded because the archived run used a non-comparable transductive independent-graph setting. GNN-CPI showed unstable training on BindingDB, with a large standard deviation across seeds.}}
\end{table}

\subsection{Token-Matched Control and Architecture-Replacement Analysis}
RQ1 is answered by comparing the full model with a token-matched standard-GCN control on the BioSNAP random split (Table~\ref{tab:replacement_configs}, Figure~\ref{fig:ablation_2x2}). The control shares the same atom featurizer, node count, drug mask, protein CNN, bilinear attention, decoder, and training protocol as the full model; the only difference is that it applies a standard multi-layer GCN over the whole molecular graph without explicit rooted-subgraph extraction. To further interpret the role of early pooling and fusion, we also compared the full model with three replacement configurations in which rooted molecular neighborhoods, bilinear drug--protein fusion, or both were replaced.

\begin{table}[H]
\caption{Architecture configurations used in the replacement analysis and the token-matched standard-GCN control. The drug-encoder replacement changed token granularity and the stage of molecular pooling, the fusion replacement changed only the interaction module, and the control omitted explicit subgraph extraction while retaining per-atom tokens.}
\label{tab:replacement_configs}
\footnotesize
\begin{tabular*}{\textwidth}{@{\extracolsep{\fill}}llllr@{}}
\toprule
\textbf{Configuration} & \textbf{Drug encoder} & \textbf{Drug tokens} & \textbf{Fusion} & \textbf{Params} \\
\midrule
Full SGBANDTI & Rooted $K$-hop subgraph & 290 & Masked bilinear & 1,070,342 \\
No subgraph & Whole-molecule GCN + mean pool & 1 & Masked bilinear & 1,070,342 \\
GCN tokens & Standard GCN (atom tokens) & 290 & Masked bilinear & 1,070,342 \\
No BAN & Rooted $K$-hop subgraph & 290 & Pooled concatenation & 935,937 \\
No both & Whole-molecule GCN + mean pool & 1 & Pooled concatenation & 935,937 \\
\bottomrule
\end{tabular*}
\noindent{\footnotesize{\textit{Note:} The rooted-subgraph configurations produce drug-token shape $[B, 290, 128]$ and the whole-molecule configurations $[B, 1, 128]$. The no-subgraph replacement reduced the drug representation from multiple atom-centered tokens to a single molecular token while keeping the parameter count unchanged.}}
\end{table} The token-matched standard-GCN control achieved mean AUROC $0.9051 \pm 0.0043$ and AUPRC $0.9110 \pm 0.0068$, close to the full model's $0.9062 \pm 0.0019$ and $0.9132 \pm 0.0043$. Paired bootstrap over the five seeds gave a mean $\Delta$AUROC of $+0.0011$ (full minus control): the 95\% confidence interval crossed zero in four of five seeds, and the full model was not consistently ranked above the control. For $\Delta$AUPRC (mean $+0.0021$), most per-seed intervals excluded zero but the sign was not consistent across seeds. Overall, no stable incremental gain from explicit rooted-subgraph extraction was observed on this split.

As supporting evidence for the role of early pooling and fusion, the three replacement configurations degraded more clearly. Replacing rooted subgraphs with a pooled whole-molecule GCN reduced AUROC from $0.9062$ to $0.8757$ and AUPRC from $0.9132$ to $0.8782$; replacing bilinear attention with pooled concatenation reduced AUROC to $0.8777$ and AUPRC to $0.8759$; and replacing both reduced AUROC to $0.8598$ and AUPRC to $0.8551$. Because the no-subgraph replacement simultaneously changes token granularity, the stage of molecular pooling, and the granularity of the interaction, these replacement comparisons describe configuration-level differences rather than isolated component ablations. The contrast between the near-matching token control and the clearly degraded pooled variants indicates that the advantage of the full configuration over pooled whole-molecule variants was mainly attributable to retaining per-atom tokens until fusion, rather than to the rooted-subgraph extraction itself.

\begin{figure}[H]
\centering
\includegraphics[width=\textwidth]{figure_ablation_2x2.png}
\caption{Architecture-replacement analysis on the BioSNAP random split. The figure compares AUROC and AUPRC across the full model, three replacement configurations, and a token-matched standard-GCN control; error bars denote the sample standard deviation ($n = 5$ training seeds). The GCN-token control performs close to the full model, whereas the pooled and concatenation variants degrade more clearly.}
\label{fig:ablation_2x2}
\end{figure}

\subsection{Generalization to Unseen Drugs and Targets}
Entity-disjoint splits were used to separate held-out-pair performance from the harder question of whether the model transfers to unseen drugs or unseen proteins. The unseen-drug split was constructed with RDKit-canonical SMILES identities, making the held-out drug setting disjoint at the canonical-SMILES level rather than only at the raw-string level. Under this setting, SGBANDTI obtained the numerically highest mean AUROC on the current frozen split ($0.8794 \pm 0.0019$), slightly ahead of DrugBAN ($0.8750 \pm 0.0032$), and the highest mean AUPRC ($0.8821 \pm 0.0026$; Figure~\ref{fig:biosnap_split_boxplots}). These results indicate the highest average drug-disjoint ranking on this single frozen split, rather than a general claim of superior cold-start generalization (Table~\ref{tab:coldstart}).

The unseen-target split was substantially harder. SGBANDTI reached an AUROC of $0.6345 \pm 0.0182$ and an AUPRC of $0.6124 \pm 0.0212$, below RF, MGNDTI, MolTrans, DrugBAN, iNGNN-DTI, and GNN-CPI in AUROC. The validation-selected threshold also transferred poorly to the unseen-target test distribution, resulting in degraded threshold-dependent performance. For this reason, AUROC and AUPRC are used as the primary evidence for unseen-target evaluation, and F1, sensitivity, specificity, and accuracy are not used to claim target-disjoint performance.

The weak unseen-target result indicates limited target-side transfer under this split, although the current experiment does not isolate the contribution of the target encoder from partition-specific distribution shift. The split-level results define SGBANDTI's strongest generalization evidence as drug-disjoint rather than protein-disjoint. Paired bootstrap gave a mean $\Delta$AUROC of $+0.0043$ for unseen-drug, with 95\% confidence intervals crossing zero in most seeds, and $-0.0215$ for unseen-target, where DrugBAN was significantly better in most seeds.

\begin{figure}[H]
\centering
\includegraphics[width=\textwidth]{figure_biosnap_split_boxplots.png}
\caption{BioSNAP split-level AUROC and AUPRC across random, canonical-SMILES unseen-drug, and protein-disjoint unseen-target settings. Diamonds mark the five-seed mean, black error bars denote the sample standard deviation ($n = 5$ training seeds), and colored points show the individual seed values.}
\label{fig:biosnap_split_boxplots}
\end{figure}

\begin{table}[H]
\caption{Cold-start generalization results on BioSNAP. Values are mean $\pm$ sample standard deviation over five seeds for the unseen-drug (canonical-SMILES-disjoint) and unseen-target (exact-protein-sequence-disjoint) splits. The best value per column is shown in bold.}
\label{tab:coldstart}
\scriptsize
\setlength{\tabcolsep}{4pt}
\renewcommand{\arraystretch}{1.05}
\begin{tabular*}{\textwidth}{@{\extracolsep{\fill}}lcccc@{}}
\toprule
\textbf{Model} & \textbf{Unseen-drug AUROC} & \textbf{Unseen-drug AUPRC} & \textbf{Unseen-target AUROC} & \textbf{Unseen-target AUPRC} \\
\midrule
SGBANDTI & \textbf{0.8794 $\pm$ 0.0019} & \textbf{0.8821 $\pm$ 0.0026} & 0.6345 $\pm$ 0.0182 & 0.6124 $\pm$ 0.0212 \\
DrugBAN & 0.8750 $\pm$ 0.0032 & 0.8806 $\pm$ 0.0052 & 0.6560 $\pm$ 0.0170 & 0.6333 $\pm$ 0.0150 \\
MGNDTI & 0.8589 $\pm$ 0.0057 & 0.8675 $\pm$ 0.0038 & 0.6910 $\pm$ 0.0105 & 0.6836 $\pm$ 0.0168 \\
RF & 0.8493 $\pm$ 0.0006 & 0.8724 $\pm$ 0.0005 & \textbf{0.6979 $\pm$ 0.0110} & \textbf{0.6867 $\pm$ 0.0068} \\
TransformerCPI & 0.8460 $\pm$ 0.0113 & 0.8551 $\pm$ 0.0061 & 0.6160 $\pm$ 0.0319 & 0.6032 $\pm$ 0.0216 \\
iNGNN-DTI & 0.8417 $\pm$ 0.0057 & 0.8510 $\pm$ 0.0023 & 0.6526 $\pm$ 0.0090 & 0.6421 $\pm$ 0.0075 \\
MolTrans & 0.8407 $\pm$ 0.0082 & 0.8473 $\pm$ 0.0057 & 0.6820 $\pm$ 0.0166 & 0.6788 $\pm$ 0.0193 \\
GNN-CPI & 0.6797 $\pm$ 0.0897 & 0.6795 $\pm$ 0.0933 & 0.6501 $\pm$ 0.0030 & 0.6526 $\pm$ 0.0017 \\
\bottomrule
\end{tabular*}
\noindent{\footnotesize{\textit{Note:} On the unseen-target split, the validation-selected threshold transferred poorly for several neural models; AUROC and AUPRC remain the primary ranking evidence.}}
\end{table}

\section{Discussion}
The main finding is that explicit rooted-subgraph extraction did not show a stable ranking gain over a token-matched standard-GCN control on the BioSNAP random split: the control reached mean AUROC $0.9051 \pm 0.0043$ and AUPRC $0.9110 \pm 0.0068$, close to the full model's $0.9062 \pm 0.0019$ and $0.9132 \pm 0.0043$, with paired-bootstrap differences inconsistent in sign across seeds. The experiments instead support delayed molecular pooling and token-level cross-modal fusion: compressing the drug representation to a single molecule-level token before fusion, or replacing bilinear attention with concatenation, clearly degraded ranking performance, whereas a token-matched control that simply omitted explicit subgraph extraction while retaining per-atom tokens performed close to the full model. This pattern indicates that the advantage of the full configuration over pooled variants mainly reflects retaining per-atom tokens until fusion, rather than the rooted-subgraph construction itself. Under random-pair splits, SGBANDTI performed close to DrugBAN; on the canonical-SMILES unseen-drug split it obtained the highest mean AUROC and AUPRC, and on the exact-protein-sequence target-disjoint split performance was limited. The protein branch supplies position-dependent sequence features, and the bilinear module aggregates latent associations between the two feature sets.

SGBANDTI is a focused molecular-graph/protein-sequence model, not a universal replacement for existing DTI architectures. GNN-CPI and GraphDTA combine molecular graphs with convolutional protein representations \citep{tsubaki2019,nguyen2021graphdta}; MolTrans models substructure-level drug--protein interactions \citep{huang2021moltrans}; DrugBAN uses low-rank bilinear attention with conventional molecular-graph and protein-sequence features \citep{bai2023drugban}; and iNGNN-DTI applies nested encoding to drug and structure-derived target graphs \citep{sun2024ingnndti}. DrugBAN already feeds standard GCN node features into bilinear attention, so the distinguishing element of SGBANDTI relative to DrugBAN is explicit rooted-subgraph extraction and masking; the token-matched control indicates that this specific difference did not yield a measurable ranking gain on the BioSNAP random split. Against this background, the contribution of SGBANDTI is narrower: it evaluates whether rooted drug-side neighborhoods add value when no protein structure, binding-site annotation, or contact supervision is supplied.

Random interaction splits can place the same drug or protein in both training and test sets and therefore cannot establish unseen-entity generalization. The processed BindingDB and BioSNAP files provide binary benchmark labels, but the released files do not fully reconstruct the original data-generation pipeline; full traceability details are consolidated in the Data Availability statement. Consequently, 0-labeled pairs should be interpreted as benchmark negatives or unobserved pairs, not as experimentally verified non-binding pairs. All repeated-seed statistics were computed on one frozen data partition, so they describe training stochasticity rather than split uncertainty. Protein sequences were truncated at 1200 residues, unrecognized residues shared the zero index with padding, and padding may still influence batch-normalization statistics before masked fusion.

The most important empirical limitation is target-disjoint generalization. SGBANDTI ranked numerically first by mean AUROC and AUPRC in the unseen-drug split but ranked low in the unseen-target split, where the validation-selected threshold transferred poorly and threshold-dependent performance degraded. Drug-side rooted subgraphs did not compensate for limited protein-side extrapolation under this evaluated setting. These explanations are hypotheses rather than directly tested mechanisms. This asymmetry is consistent with the architecture: the main representational novelty lies on the drug side, whereas the protein branch is a supervised sequence CNN with limited external biological prior. When a test protein has no training-set occurrence, the model must infer target similarity from local amino-acid patterns alone. This signal may be insufficient for remote protein-family relationships, domain-level function, binding-site compatibility, or target-class shifts that are not represented in the training labels. Stronger protein representations, protein-family-aware evaluation, pretrained sequence encoders, or structure-aware target features may therefore be needed before the model can be used for robust target cold-start prioritization. In addition, the bilinear association maps are unnormalized latent scores and should not be read as binding contacts or interaction energies.

Future work could address target-disjoint evaluation, repeated alternative partitions, and prospective or temporally separated biological validation.

\section{Conclusions}
Whether explicit rooted-subgraph extraction improves DTI ranking beyond a standard GCN that retains the same atom-level tokens until bilinear fusion remains an open question for the settings evaluated here. SGBANDTI, which combines atom-centered molecular-subgraph encoding, sequence-based protein convolution, and masked low-rank bilinear feature interaction, was competitive on BioSNAP and BindingDB random-pair splits and showed a small numerical advantage in mean ranking metrics on the canonical-SMILES unseen-drug split. A token-matched standard-GCN control achieved similar mean performance with seed-dependent differences (mean AUROC $0.9051 \pm 0.0043$ versus $0.9062 \pm 0.0019$ for the full model), indicating that the improvement over pooled whole-molecule variants was mainly associated with retaining atom-level tokens until fusion rather than with explicit rooted-subgraph extraction. SGBANDTI was competitive on random-pair splits but exhibited limited transfer under the evaluated exact-protein-sequence-disjoint protocol, which remains the main limitation.

\authorcontributions{Conceptualization, W.L. and X.Z.; methodology, W.L., Q.L. and M.L.; software, W.L. and Q.L.; validation, Q.L., Y.F. and M.L.; formal analysis, W.L. and J.X.; investigation, W.L.; resources, X.Z. and Y.Y.; data curation, W.L. and Q.L.; visualization, W.L. and M.L.; writing---original draft preparation, W.L. and Q.L.; writing---review and editing, X.Z., J.X. and Y.Y.; supervision, X.Z.; project administration, X.Z. and Y.Y.; funding acquisition, W.L., J.X. and Y.Y. All authors have read and agreed to the published version of the manuscript.}

\funding{This work was supported in part by the National Natural Science Foundation of China (62472051), Hunan Provincial Natural Science Foundation of China (2026JJ80253), Hunan Provincial Science and Technology Department 2024 Joint Natural Science Foundation for Medical and Health Industry (2024JJ9473), and the Scientific Research Fund of Hunan Provincial Education Department (24C0556).}

\institutionalreview{Not applicable.}

\informedconsent{Not applicable.}

\dataavailability{The source code and processed benchmark files supporting the reported results are available at \url{https://github.com/BMJ520tc/-SGBANDTI-paper}. The source resources are BindingDB \citep{liu2007bindingdb} and BioSNAP \citep{zitnik2018}, which are available from the providers cited in the reference list; the BioSNAP entry cites the Stanford BioSNAP dataset page (accessed 31 August 2026). The processed files were taken from the public code repositories of DrugBAN \citep{bai2023drugban} and iNGNN-DTI \citep{sun2024ingnndti}, with sample counts matching the frozen split records in this repository. The released processed files contain SMILES strings, protein sequences, and binary benchmark labels. Preprocessing followed a reproducible pipeline: drug SMILES were canonicalized with RDKit, duplicate pairs were removed (seven pairs in BioSNAP; conflicting labels retained the positive label), and each split was generated with \texttt{build\_splits.py}. Frozen split hashes and sample counts are recorded in the machine-readable split metadata accompanying this repository. Reproduction is supported from these processed CSV files and the frozen split records. The released files do not fully reconstruct the original BindingDB and BioSNAP dataset-generation procedures, specifically the BindingDB release date, affinity-to-binary conversion rule, BindingDB negative-pair construction, BioSNAP snapshot date, and BioSNAP unlabeled-pair sampling rule. Accordingly, 0-labeled pairs are treated as benchmark negative or unobserved pairs rather than as experimentally verified non-binding pairs.}

\supplementary{The following supporting information will be submitted with the manuscript: frozen split metadata; processed train, validation, and test CSV files; per-seed metric and prediction files; figure-generation scripts; duplicate and label-conflict resolution logs; source-file checksums; a baseline protocol table listing each comparator's code source, input modalities, seeds, test-set sizes, and AUPRC definition; and a data-provenance note specifying, where available, the BindingDB release and download date, affinity-to-binary thresholding rule, BindingDB negative-sampling procedure, BioSNAP snapshot, and BioSNAP unlabeled-pair sampling procedure.}

\acknowledgments{The authors thank the maintainers of BindingDB and BioSNAP for making high-quality biomedical resources publicly available. The authors also acknowledge the developers of the open-source software tools used in this work, which supported model development, evaluation, and visualization. The use of generative AI during the preparation of this manuscript is disclosed in the Materials and Methods section.}

\conflictsofinterest{The authors declare no conflicts of interest. The funders had no role in the design of the study; in the collection, analyses, or interpretation of data; in the writing of the manuscript; or in the decision to publish the results.}

\abbreviations{Abbreviations}{%
\begin{tabular}{@{}ll}
DTI & Drug--target interaction \\
GNN & Graph neural network \\
GCN & Graph convolutional network \\
CNN & Convolutional neural network \\
MLP & Multilayer perceptron \\
BAN & Bilinear attention network \\
BCE & Binary cross-entropy \\
RF & Random Forest \\
AUROC & Area under the receiver operating characteristic curve \\
AUPRC & Area under the precision--recall curve
\end{tabular}}

\begin{adjustwidth}{-\extralength}{0cm}
\reftitle{References}
\bibliography{sample}
\end{adjustwidth}

\end{document}
